\documentclass[12pt]{article}
\usepackage[utf8]{inputenc}
\usepackage{amsmath, amssymb, amsthm}
\usepackage[margin=1in]{geometry}
\usepackage{enumitem}
\usepackage{xcolor}

% ── Theorem-style environments ──
\newtheorem{definition}{Definition}
\newtheorem{theorem}{Theorem}
\newtheorem{example}{Example}

% ── Custom commands ──
\newcommand{\vect}[1]{\vec{#1}}
\newcommand{\uvec}[1]{\hat{#1}}
\newcommand{\grad}{\nabla}
\newcommand{\norm}[1]{\left\| #1 \right\|}
\newcommand{\inner}[2]{\left\langle #1,\, #2 \right\rangle}

\title{\textbf{Midterm Mock Test Solutions} \\ \large MTH253 -- Calculus III}
\author{}
\date{}

\begin{document}
\maketitle

%% ================================================================
%%  QUESTION 2: Directional Derivatives
%% ================================================================
\section{Question 2 -- Directional Derivatives}

\subsection{Key Formulas}

\begin{definition}[Directional Derivative]
The \textbf{directional derivative} of $f(x,y)$ in the direction of a
\textbf{unit vector} $\vect{u} = \inner{a}{b}$ is
\[
  D_{\vect{u}}\, f(x,y)
    = f_x(x,y)\,a + f_y(x,y)\,b
    = \grad f(x,y) \cdot \vect{u}.
\]
\end{definition}

\medskip

\noindent\textbf{Two common situations:}

\begin{enumerate}[label=(\roman*)]
  \item \textbf{Direction given by an angle $\theta$:}
        Set $\vect{u} = \inner{\cos\theta}{\sin\theta}$, so
        \[
          D_{\vect{u}}\, f(x,y)
            = f_x(x,y)\cos\theta + f_y(x,y)\sin\theta.
        \]

  \item \textbf{Direction given by a (non-unit) vector $\vect{v}$:}
        You \textbf{must} first normalise:
        \[
          \vect{u} = \frac{\vect{v}}{\norm{\vect{v}}},
        \]
        then apply the formula above.
\end{enumerate}

\subsection{General Solution Method}

Given $f(x,y)$, a point $P=(x_0,y_0)$, and a direction (angle or vector):

\begin{enumerate}
  \item \textbf{Compute the partial derivatives} $f_x$ and $f_y$.
  \item \textbf{Evaluate at the point:} $f_x(x_0,y_0)$ and $f_y(x_0,y_0)$.
  \item \textbf{Form the gradient:}
        $\grad f(x_0,y_0) = \inner{f_x(x_0,y_0)}{f_y(x_0,y_0)}$.
  \item \textbf{Obtain the unit direction vector $\vect{u}$:}
        \begin{itemize}
          \item If an angle $\theta$ is given:
                $\vect{u}=\inner{\cos\theta}{\sin\theta}$.
          \item If a vector $\vect{v}$ is given:
                $\vect{u}=\dfrac{\vect{v}}{\norm{\vect{v}}}$.
        \end{itemize}
  \item \textbf{Dot product:}
        $D_{\vect{u}}\,f(x_0,y_0) = \grad f(x_0,y_0)\cdot\vect{u}$.
\end{enumerate}

\subsection{Worked Example 1 -- Direction Given by an Angle}

\begin{example}
Find the directional derivative of
$f(x,y) = x^2 y^3 - 4y$
at the point $(2,-1)$ in the direction $\theta = \pi/6$.
\end{example}

\textbf{Solution.}

\medskip
\noindent\textbf{Step 1.} Partial derivatives:
\[
  f_x = 2xy^3, \qquad f_y = 3x^2 y^2 - 4.
\]

\noindent\textbf{Step 2.} Evaluate at $(2,-1)$:
\[
  f_x(2,-1) = 2(2)(-1)^3 = -4, \qquad
  f_y(2,-1) = 3(4)(1) - 4 = 8.
\]

\noindent\textbf{Step 3.} Gradient:
\[
  \grad f(2,-1) = \inner{-4}{8}.
\]

\noindent\textbf{Step 4.} Unit vector from angle $\theta=\pi/6$:
\[
  \vect{u}
    = \inner{\cos\frac{\pi}{6}}{\sin\frac{\pi}{6}}
    = \inner{\frac{\sqrt{3}}{2}}{\frac{1}{2}}.
\]

\noindent\textbf{Step 5.} Directional derivative:
\[
  D_{\vect{u}}\,f(2,-1)
    = \inner{-4}{8} \cdot \inner{\frac{\sqrt{3}}{2}}{\frac{1}{2}}
    = -4\cdot\frac{\sqrt{3}}{2} + 8\cdot\frac{1}{2}
    = -2\sqrt{3} + 4.
\]

\[
  \boxed{D_{\vect{u}}\,f(2,-1) = 4 - 2\sqrt{3} \approx 0.536}
\]

\subsection{Worked Example 2 -- Direction Given by a Vector}

\begin{example}
Find the directional derivative of
$f(x,y) = e^x \sin y$
at $(0,\,\pi/3)$ in the direction of $\vect{v}=\inner{3}{4}$.
\end{example}

\textbf{Solution.}

\medskip
\noindent\textbf{Step 1.} Partial derivatives:
\[
  f_x = e^x \sin y, \qquad f_y = e^x \cos y.
\]

\noindent\textbf{Step 2.} Evaluate at $(0,\,\pi/3)$:
\[
  f_x\!\left(0,\tfrac{\pi}{3}\right) = e^0 \sin\frac{\pi}{3}
    = \frac{\sqrt{3}}{2}, \qquad
  f_y\!\left(0,\tfrac{\pi}{3}\right) = e^0 \cos\frac{\pi}{3}
    = \frac{1}{2}.
\]

\noindent\textbf{Step 3.} Gradient:
\[
  \grad f\!\left(0,\tfrac{\pi}{3}\right)
    = \inner{\frac{\sqrt{3}}{2}}{\frac{1}{2}}.
\]

\noindent\textbf{Step 4.} Normalise $\vect{v}=\inner{3}{4}$:
\[
  \norm{\vect{v}} = \sqrt{9+16} = 5, \qquad
  \vect{u} = \inner{\frac{3}{5}}{\frac{4}{5}}.
\]

\noindent\textbf{Step 5.} Directional derivative:
\[
  D_{\vect{u}}\,f\!\left(0,\tfrac{\pi}{3}\right)
    = \inner{\frac{\sqrt{3}}{2}}{\frac{1}{2}}
      \cdot \inner{\frac{3}{5}}{\frac{4}{5}}
    = \frac{3\sqrt{3}}{10} + \frac{4}{10}
    = \frac{3\sqrt{3}+4}{10}.
\]

\[
  \boxed{D_{\vect{u}}\,f\!\left(0,\tfrac{\pi}{3}\right)
    = \frac{3\sqrt{3}+4}{10} \approx 0.920}
\]

%% ================================================================
%%  SECTION 5: QUESTION 3 -- Extreme Values on a Disk
%% ================================================================
\newpage
\section{Question 3 -- Extreme Values on a Closed Disk}

\begin{example}
Find the extreme values of the function
\[
  f(x,y) = x^2 + y^2 + 4x - 4y
\]
on the disk $x^2 + y^2 \le 9$.
\end{example}

\textbf{Solution.}

\medskip
Since the disk $D = \{(x,y) : x^2+y^2 \le 9\}$ is a \textbf{closed and bounded}
region, by the Extreme Value Theorem $f$ attains both an absolute maximum and an
absolute minimum on $D$.  We must check:
\begin{enumerate}[label=(\alph*)]
  \item Interior critical points (where $\grad f = \vect{0}$).
  \item The boundary $x^2+y^2 = 9$.
\end{enumerate}

%% ── Part (a): Interior ──
\subsection*{Part (a): Interior critical points}

Compute the partial derivatives and set them to zero:
\[
  f_x = 2x + 4 = 0 \;\Longrightarrow\; x = -2, \qquad
  f_y = 2y - 4 = 0 \;\Longrightarrow\; y = 2.
\]

The only critical point is $(-2, 2)$.  Check whether it lies inside the disk:
\[
  (-2)^2 + 2^2 = 4 + 4 = 8 \le 9 \quad \checkmark
\]

So $(-2,2)$ is in the disk. Evaluate $f$ there:
\[
  f(-2,2) = 4 + 4 + 4(-2) - 4(2) = 8 - 8 - 8 = \mathbf{-8}.
\]

%% ── Part (b): Boundary ──
\subsection*{Part (b): Boundary $x^2 + y^2 = 9$}

Parametrise the circle of radius $3$:
\[
  x = 3\cos \theta, \qquad y = 3\sin \theta, \qquad 0 \le \theta < 2\pi.
\]

Substitute into $f$:
\begin{align*}
  g(\theta) &= (3\cos \theta)^2 + (3\sin \theta)^2 + 4(3\cos \theta) - 4(3\sin \theta) \\
       &= 9\cos^2 \theta + 9\sin^2 \theta + 12\cos \theta - 12\sin \theta \\
       &= 9 + 12\cos \theta - 12\sin \theta.
\end{align*}

To find extrema of $g(\theta)$, set $g'(\theta) = 0$:
\[
  g'(\theta) = -12\sin \theta - 12\cos \theta = 0
  \;\Longrightarrow\;
  \sin \theta = -\cos \theta
  \;\Longrightarrow\;
  \tan \theta = -1.
\]

This gives two solutions in $[0, 2\pi)$:
\[
  \theta_1 = \frac{3\pi}{4}, \qquad \theta_2 = \frac{7\pi}{4}.
\]

Evaluate $g$ at each:

\medskip
\noindent\textbf{At $\theta_1 = \dfrac{3\pi}{4}$:}
\[
  \cos\frac{3\pi}{4} = -\frac{\sqrt{2}}{2}, \qquad
  \sin\frac{3\pi}{4} = \frac{\sqrt{2}}{2}.
\]
\[
  g\!\left(\tfrac{3\pi}{4}\right)
    = 9 + 12\!\left(-\tfrac{\sqrt{2}}{2}\right) - 12\!\left(\tfrac{\sqrt{2}}{2}\right)
    = 9 - 6\sqrt{2} - 6\sqrt{2}
    = 9 - 12\sqrt{2}.
\]

\noindent\textbf{At $\theta_2 = \dfrac{7\pi}{4}$:}
\[
  \cos\frac{7\pi}{4} = \frac{\sqrt{2}}{2}, \qquad
  \sin\frac{7\pi}{4} = -\frac{\sqrt{2}}{2}.
\]
\[
  g\!\left(\tfrac{7\pi}{4}\right)
    = 9 + 12\!\left(\tfrac{\sqrt{2}}{2}\right) - 12\!\left(-\tfrac{\sqrt{2}}{2}\right)
    = 9 + 6\sqrt{2} + 6\sqrt{2}
    = 9 + 12\sqrt{2}.
\]

The corresponding points on the boundary are:
\[
  \theta_1:\quad \left(-\tfrac{3\sqrt{2}}{2},\;\tfrac{3\sqrt{2}}{2}\right),
  \qquad
  \theta_2:\quad \left(\tfrac{3\sqrt{2}}{2},\;-\tfrac{3\sqrt{2}}{2}\right).
\]

%% ── Conclusion ──
\subsection*{Conclusion}

Compare all candidate values:

\medskip
\begin{center}
\renewcommand{\arraystretch}{1.4}
\begin{tabular}{|c|c|c|}
\hline
\textbf{Point} & \textbf{Type} & $\boldsymbol{f(x,y)}$ \\
\hline
$(-2,\;2)$ & Interior critical pt. & $-8$ \\
\hline
$\left(-\frac{3\sqrt{2}}{2},\;\frac{3\sqrt{2}}{2}\right)$
  & Boundary & $9 - 12\sqrt{2} \approx -7.97$ \\
\hline
$\left(\frac{3\sqrt{2}}{2},\;-\frac{3\sqrt{2}}{2}\right)$
  & Boundary & $9 + 12\sqrt{2} \approx 25.97$ \\
\hline
\end{tabular}
\end{center}

\bigskip
\[
  \boxed{\text{Absolute minimum} = -8
    \;\text{ at }\; (-2,\,2)}
\]
\[
  \boxed{\text{Absolute maximum} = 9 + 12\sqrt{2}
    \;\text{ at }\;
    \left(\tfrac{3\sqrt{2}}{2},\;-\tfrac{3\sqrt{2}}{2}\right)}
\]

%% ================================================================
%%  SECTION 6: QUESTION 4 -- Volume via Double Integral
%% ================================================================
\newpage
\section{Question 4 -- Volume via Double Integral}

\begin{example}
Find the volume of the solid under the surface
\[
  z = 1 + x^2 y^2
\]
and above the region enclosed by $x = y^2$ and $x = 4$.
\end{example}

\textbf{Solution.}

\medskip
The volume is given by the double integral
\[
  V = \iint_R \bigl(1 + x^2 y^2\bigr)\, dA,
\]
where $R$ is the region in the $xy$-plane bounded by the parabola $x = y^2$
and the vertical line $x = 4$.

\subsection*{Step 1: Sketch the region and determine limits}

The curves $x = y^2$ and $x = 4$ intersect when
\[
  y^2 = 4 \;\Longrightarrow\; y = \pm 2.
\]

For a fixed $y \in [-2,\,2]$, $x$ runs from the parabola to the line:
\[
  y^2 \le x \le 4.
\]

So the integral becomes
\[
  V = \int_{-2}^{2} \int_{y^2}^{4} \bigl(1 + x^2 y^2\bigr)\, dx\, dy.
\]

\subsection*{Step 2: Evaluate the inner integral (with respect to $x$)}

\begin{align*}
  \int_{y^2}^{4} \bigl(1 + x^2 y^2\bigr)\, dx
    &= \left[ x + \frac{x^3}{3}\, y^2 \right]_{x=y^2}^{x=4} \\
    &= \left(4 + \frac{64}{3}\, y^2\right)
       - \left(y^2 + \frac{y^6}{3}\, y^2\right) \\
    &= 4 + \frac{64}{3}\, y^2 - y^2 - \frac{y^8}{3} \\
    &= 4 + \frac{64 y^2 - 3y^2}{3} - \frac{y^8}{3} \\
    &= 4 + \frac{61 y^2}{3} - \frac{y^8}{3}.
\end{align*}

\subsection*{Step 3: Evaluate the outer integral (with respect to $y$)}

Since the integrand $4 + \frac{61 y^2}{3} - \frac{y^8}{3}$ is an
\textbf{even function} of $y$ (only even powers), we can use symmetry:
\[
  V = 2 \int_{0}^{2} \left(4 + \frac{61 y^2}{3} - \frac{y^8}{3}\right) dy.
\]

Evaluate each term:
\begin{align*}
  \int_{0}^{2} 4\, dy &= 4 \cdot 2 = 8, \\[6pt]
  \int_{0}^{2} \frac{61 y^2}{3}\, dy
    &= \frac{61}{3} \cdot \frac{y^3}{3}\bigg|_0^2
     = \frac{61}{3} \cdot \frac{8}{3}
     = \frac{488}{9}, \\[6pt]
  \int_{0}^{2} \frac{y^8}{3}\, dy
    &= \frac{1}{3} \cdot \frac{y^9}{9}\bigg|_0^2
     = \frac{1}{3} \cdot \frac{512}{9}
     = \frac{512}{27}.
\end{align*}

Combine:
\begin{align*}
  \frac{V}{2}
    &= 8 + \frac{488}{9} - \frac{512}{27} \\
    &= \frac{216}{27} + \frac{1464}{27} - \frac{512}{27}
     \quad\text{(common denominator 27)} \\
    &= \frac{216 + 1464 - 512}{27} \\
    &= \frac{1168}{27}.
\end{align*}

Therefore:
\[
  V = 2 \cdot \frac{1168}{27} = \frac{2336}{27}.
\]

\[
  \boxed{V = \dfrac{2336}{27} \approx 86.52}
\]

\end{document}
